\section{Parámetros adicionales de RMSNorm}

\subsection{Estructura de RMSNorm}

RMSNorm (Root Mean Square Normalization) es el método de normalización introducido por LLaMA 2, más simple que LayerNorm: solo hace escalamiento, sin centrado de la media:

$$
\text{RMSNorm}(x) = \frac{x}{\sqrt{\frac{1}{d}\sum_{i=1}^d x_i^2 + \epsilon}} \odot \gamma
$$

donde $\gamma \in \mathbb{R}^{d_{\text{model}}}$ es el parámetro de escalamiento entrenable.

\subsection{Distribución de RMSNorm}

Cada Transformer Layer contiene \textbf{2} RMSNorm:
\begin{enumerate}
  \item la Norm antes de GQA Attention
  \item la Norm antes de FFN/MoE
\end{enumerate}

Sumando el \textbf{1} Final Norm al final del modelo, en total:

$$
P_{\text{RMSNorm}} = (2 \times L + 1) \times d_{\text{model}}
$$

Para Dense 32B: $(2 \times 64 + 1) \times 5120 = 660{,}480$ parámetros. Aunque es despreciable frente a los 32B, no debe omitirse conceptualmente.

\subsection{Resumen de la sección}

RMSNorm es «estándar» en cada capa: su número de parámetros es pequeño pero indispensable. Garantiza la estabilidad numérica de la entrada de cada capa y es uno de los componentes clave para entrenar Transformers profundos.

